% ---------------------------------------------------------------------------
% Chapter: neural_toolkit
% ---------------------------------------------------------------------------
\chapter{neural\_toolkit: Networks and Tabular Tools}\label{chap:neural}
% Long monospaced identifiers leave little room for line breaking; allow some
% extra inter-word stretch in this chapter (the group keeps it local).
\begingroup\setlength{\emergencystretch}{3em}% chapter-local

\pkg{toolkit.neural\_toolkit} provides configurable PyTorch modules for
reinforcement learning: policy, value and Q-networks with MLP, convolutional,
recurrent and transformer trunks, encoders and decoders (including a
variational autoencoder), factories that build them by name, a collection of
utilities for initialisation, optimisation and checkpoints, and NumPy-based
tools for tabular methods. It does not implement learning algorithms; the
modules are ordinary \code{torch.nn.Module} objects meant to be combined with
your own training code.

The subpackage requires PyTorch (\code{pip install ".[torch]"}). Importing it
without PyTorch raises an \code{ImportError} that explains how to install it;
the rest of \pkg{toolkit} does not need PyTorch.

\section{Conventions}\label{sec:neural-conventions}

\begin{itemize}
  \item \emph{Batch first.} Inputs have a leading batch dimension \code{B}:
    feature vectors \code{(B, D)}, images \code{(B, C, H, W)} and sequences
    \code{(B, T, D)}.
  \item \emph{One MLP convention.} Every fully connected stack in the
    toolkit consists of hidden blocks
    \code{Linear -> [LayerNorm] -> activation -> [Dropout]} followed by a
    linear output layer without activation.
  \item \emph{Raw outputs.} Policies return unnormalised outputs of size
    \code{output\_dim}: logits for a categorical policy, or distribution
    parameters (for example means) for a continuous one. Turning them into a
    distribution is left to the algorithm.
  \item \emph{Devices.} Every constructor accepts \code{device} (default
    \code{"cpu"}) and moves the module there at the end of construction;
    \code{device=None} leaves it where PyTorch created it. Inputs must be on
    the same device.
  \item \emph{Validation.} Invalid sizes, dropout rates, activation names or
    cell types raise \code{ValueError} or \code{TypeError} at construction,
    with a message that lists the valid options.
\end{itemize}

\section{Building blocks}\label{sec:neural-blocks}

\paragraph{Activations.}
\code{get\_activation(name)} returns a new activation module for a
case-insensitive name, or instantiates an \code{nn.Module} subclass passed
instead of a name. The valid names are available as \code{ACTIVATION\_NAMES}:

\begin{center}
  \small
  \begin{tabular}{@{}llll@{}}
    \toprule
    Name & Module & Name & Module \\
    \midrule
    \code{relu}        & \code{nn.ReLU}      & \code{gelu}     & \code{nn.GELU} \\
    \code{leaky\_relu} & \code{nn.LeakyReLU} & \code{silu}, \code{swish} & \code{nn.SiLU} \\
    \code{elu}         & \code{nn.ELU}       & \code{mish}     & \code{nn.Mish} \\
    \code{tanh}        & \code{nn.Tanh}      & \code{softplus} & \code{nn.Softplus} \\
    \code{sigmoid}     & \code{nn.Sigmoid}   & \code{identity} & \code{nn.Identity} \\
    \bottomrule
  \end{tabular}
\end{center}

\paragraph{MLPs.}
\code{build\_mlp} returns an \code{nn.Sequential}. Without \code{output\_dim}
only the hidden stack is built; its output width is \code{hidden\_dims[-1]}.

% norun
\begin{pycode}
build_mlp(input_dim, hidden_dims, output_dim=None, activation="relu", dropout=0.0,
          layer_norm=False, output_activation=None, bias=True) -> nn.Sequential
\end{pycode}

\begin{pycode}
import torch
from toolkit.neural_toolkit import build_mlp, get_activation

mlp = build_mlp(4, [32, 32], output_dim=2, activation="tanh", layer_norm=True)
print(mlp(torch.zeros(8, 4)).shape)       # torch.Size([8, 2])
print(get_activation("swish"))            # SiLU()
\end{pycode}

\paragraph{Positional encoding.}
\code{PositionalEncoding(d\_model, dropout=0.1, max\_len=5000)} adds the
sinusoidal table of Vaswani et al.\ (2017) to a batch-first sequence
\code{(B, T, d\_model)}; position $t$ receives
\begin{equation}
  \mathrm{PE}_{t,2i} = \sin\!\left(\frac{t}{10000^{2i/d}}\right), \qquad
  \mathrm{PE}_{t,2i+1} = \cos\!\left(\frac{t}{10000^{2i/d}}\right),
\end{equation}
for every batch element, followed by dropout. Sequences longer than
\code{max\_len} raise a \code{ValueError}. The table is a non-persistent
buffer and is not stored in \code{state\_dict}.

\section{Network families}\label{sec:neural-networks}

Table~\ref{tab:neural-shapes} lists all networks with their input and output
shapes. Four classes have two names: the historical spellings
\code{CNPolicyNetwork}, \code{RNPolicyNetwork}, \code{CNQNetwork} and
\code{RNQNetwork} and the aliases with \code{CNN} and \code{RNN} refer to the
same classes.

\begin{table}[htbp]
  \centering
  \small
  \caption{Networks, encoders and decoders with their shapes.
    \code{B}: batch, \code{T}: sequence length, \code{A}: \code{action\_dim},
    \code{O}: \code{output\_dim}, \code{L}: \code{latent\_dim}.}
  \label{tab:neural-shapes}
  \begin{tabularx}{\textwidth}{@{}>{\raggedright\arraybackslash}p{5.0cm}>{\raggedright\arraybackslash}p{4.1cm}>{\raggedright\arraybackslash}X@{}}
    \toprule
    Class (alias) & Input & Output \\
    \midrule
    \multicolumn{3}{@{}l}{\emph{Policy networks} (\code{input\_dim}, \code{output\_dim})} \\
    \code{MLPPolicyNetwork} & \code{(B, D)} & \code{(B, O)} \\
    \code{CNPolicyNetwork} (\code{CNNPolicyNetwork}) & \code{(B, C, H, W)} & \code{(B, O)} \\
    \code{RNPolicyNetwork} (\code{RNNPolicyNetwork}) & \code{(B, T, D)} or \code{(B, D)} & \code{((B, O), hidden)} \\
    \code{TransformerPolicyNetwork} & \code{(B, T, D)}, mask \code{(B, T)} & \code{(B, O)} \\
    \midrule
    \multicolumn{3}{@{}l}{\emph{Value networks} (\code{input\_dim}, \code{output\_dim=1})} \\
    \code{MLPValueNetwork} & \code{(B, D)} & \code{(B, O)} \\
    \code{CNNValueNetwork} & \code{(B, C, H, W)} & \code{(B, O)} \\
    \code{RNNValueNetwork} & \code{(B, T, D)} or \code{(B, D)} & \code{((B, O), hidden)} \\
    \code{TransformerValueNetwork} & \code{(B, T, D)}, mask \code{(B, T)} & \code{(B, O)} \\
    \midrule
    \multicolumn{3}{@{}l}{\emph{Q-networks} (\code{state\_dim}, \code{action\_dim}); optional integer \code{action} \code{(B,)}} \\
    \code{MLPQNetwork} & \code{(B, D)} & \code{(B, A)}, or \code{(B,)} with \code{action} \\
    \code{MLPQNetwork}, continuous critic & \code{(B, D)}, \code{(B, A)} & \code{(B, 1)}; see below \\
    \code{DuelingQNetwork} & \code{(B, D)} & \code{(B, A)}, or \code{(B,)} with \code{action} \\
    \code{CNQNetwork} (\code{CNNQNetwork}) & \code{(B, C, H, W)} & \code{(B, A)}, or \code{(B,)} with \code{action} \\
    \code{RNQNetwork} (\code{RNNQNetwork}) & \code{(B, T, D)} or \code{(B, D)} & \code{(q, hidden)}, \code{q} as above \\
    \code{TransformerQNetwork} & \code{(B, T, D)}, mask \code{(B, T)} & \code{(B, A)}, or \code{(B,)} with \code{action} \\
    \midrule
    \multicolumn{3}{@{}l}{\emph{Encoders} (\code{input\_dim} or \code{input\_channels}, \code{latent\_dim})} \\
    \code{MLPEncoder} & \code{(B, D)} & \code{(B, L)} \\
    \code{CNNEncoder} & \code{(B, C, H, W)} & \code{(B, L)} \\
    \code{RNNEncoder} & \code{(B, T, D)} or \code{(B, D)} & \code{((B, L), hidden)} \\
    \code{TransformerEncoder} & \code{(B, T, D)}, mask \code{(B, T)} & \code{(B, L)} \\
    \code{VariationalEncoder} & \code{(B, D)} & \code{(z, mu, logvar)}, each \code{(B, L)} \\
    \midrule
    \multicolumn{3}{@{}l}{\emph{Decoders} (\code{latent\_dim}, \code{output\_dim} or \code{output\_channels})} \\
    \code{MLPDecoder}, \code{VariationalDecoder} & \code{(B, L)} & \code{(B, O)} \\
    \code{CNNDecoder} & \code{(B, L)} & \code{(B, C, S, S)}, $S = \text{\code{initial\_size}} \cdot \prod \text{\code{strides}}$ \\
    \code{RNNDecoder}, \code{TransformerDecoder} & \code{(B, L)}, \code{seq\_len} & \code{(B, T, O)}, \code{T = seq\_len} or \code{max\_seq\_len} \\
    \bottomrule
  \end{tabularx}
\end{table}

\subsection{Constructor arguments}\label{sec:neural-arguments}

All classes share the arguments \code{activation="relu"},
\code{dropout=0.0}, \code{layer\_norm=False} and \code{device="cpu"}; the
transformer classes default to \code{dropout=0.1} and \code{layer\_norm=True}.
The remaining arguments depend on the trunk (Table~\ref{tab:neural-trunks}).

\begin{table}[htbp]
  \centering
  \small
  \caption{Trunk-specific constructor arguments and their defaults.}
  \label{tab:neural-trunks}
  \begin{tabularx}{\textwidth}{@{}>{\raggedright\arraybackslash}p{2.1cm}>{\raggedright\arraybackslash}X@{}}
    \toprule
    Trunk & Arguments \\
    \midrule
    MLP & \code{hidden\_dims=(256, 256)}. For \code{DuelingQNetwork} the shared
      trunk uses \code{hidden\_dims[:-1]} and each stream one hidden layer of
      width \code{hidden\_dims[-1]}. \\
    Convolutional &
      \code{input\_channels}, \code{conv\_dims=(32, 64, 128)},
      \code{fc\_dims=(256, 256)}, \code{kernel\_sizes} (int or one per layer,
      default 3), \code{strides} (default 1). Padding is
      \code{kernel\_size // 2}. The feature map is globally average-pooled, so
      any image size works. \code{CNNEncoder}: \code{conv\_dims=(32, 64, 128,
      256)}, \code{fc\_dims=(512, 256)}, strides 1 for the first layer and 2
      afterwards, \code{pooled\_size=4} (the map is pooled to $4 \times 4$). \\
    Transposed conv. &
      \code{CNNDecoder}: \code{output\_channels}, \code{fc\_dims=(512, 256)},
      \code{conv\_dims=(256, 128, 64, 32)}, \code{kernel\_sizes} (default 3),
      \code{strides} (default 2), \code{initial\_size=4},
      \code{output\_activation="sigmoid"}; the default output is
      $64 \times 64$ (property \code{output\_size}). \\
    Recurrent &
      \code{hidden\_dim=256}, \code{num\_layers=2}, \code{rnn\_type="lstm"}
      (or \code{"gru"}), \code{bidirectional=False} (\code{True} for
      \code{RNNEncoder}; not available for \code{RNNDecoder}),
      \code{fc\_dims=(256,)}; \code{RNNDecoder}: \code{max\_seq\_len=100}. \\
    Transformer &
      \code{d\_model=256} (divisible by \code{nhead}), \code{nhead=8},
      \code{num\_layers=6}, \code{dim\_feedforward=1024}, \code{dropout=0.1},
      \code{fc\_dims=(256,)}, \code{max\_len=5000};
      \code{TransformerDecoder}: \code{max\_seq\_len=100},
      \code{max\_len=None} (then \code{max(max\_seq\_len, 5000)}). \\
    \bottomrule
  \end{tabularx}
\end{table}

\code{layer\_norm=True} inserts \code{LayerNorm} after every hidden linear
layer; in convolutional trunks it additionally inserts \code{BatchNorm2d}
after every convolution.

\subsection{Architectures}\label{sec:neural-architectures}

\paragraph{Convolutional trunk.}
\code{Conv2d -> [BatchNorm2d] -> activation -> [Dropout2d]} blocks, adaptive
average pooling to \code{pooled\_size} $\times$ \code{pooled\_size} (1 for
policies, value and Q-networks), flattening and an MLP with \code{fc\_dims},
followed by the output layer. Attribute names: \code{conv\_layers},
\code{adaptive\_pool}, \code{fc\_layers}.

\paragraph{Recurrent trunk.}
A batch-first LSTM or GRU followed by an MLP head that reads the final hidden
state of the last layer; for a bidirectional RNN the forward state after the
last step and the backward state after the first step are concatenated, so
both directions have seen the whole sequence. \code{forward(x, hidden=None)}
accepts a sequence or a single step \code{(B, D)} and returns the new hidden
state (\code{h\_n} for a GRU, \code{(h\_n, c\_n)} for an LSTM), which can be
passed to the next call:

\begin{pycode}
from toolkit.neural_toolkit import RNNPolicyNetwork

policy = RNNPolicyNetwork(input_dim=5, output_dim=3, hidden_dim=32, rnn_type="gru")
hidden = None
for t in range(4):                         # one observation per call
    logits, hidden = policy(torch.randn(8, 5), hidden)
print(logits.shape, hidden.shape)          # (8, 3) and (num_layers, 8, 32)
\end{pycode}

\paragraph{Transformer trunk.}
A linear projection to \code{d\_model}, the positional encoding, a batch-first
\code{nn.TransformerEncoder}, mean pooling over the valid positions and an MLP
head. The optional \code{mask} is a boolean \code{(B, T)} tensor in which
\code{True} marks padding (PyTorch's \code{src\_key\_padding\_mask}
convention); padded positions are excluded from attention and pooling.
Any sequence length up to \code{max\_len} is accepted.

\begin{pycode}
from toolkit.neural_toolkit import TransformerPolicyNetwork

policy = TransformerPolicyNetwork(input_dim=6, output_dim=4, d_model=64, nhead=4,
                                  num_layers=2, dim_feedforward=128)
x = torch.randn(2, 10, 6)                          # two sequences of length 10
mask = torch.zeros(2, 10, dtype=torch.bool)
mask[1, 7:] = True                                 # second sequence has 7 steps
print(policy(x, mask).shape)                       # torch.Size([2, 4])
\end{pycode}

\paragraph{Q-networks.}
Discrete Q-networks return one value per action, \code{(B, A)}. Passing
integer action indices of shape \code{(B,)} or \code{(B, 1)} as second
argument returns the selected values $Q(s, a)$ with shape \code{(B,)}, which
is what temporal-difference losses need. \code{MLPQNetwork} with
\code{continuous\_action=True} is a critic for continuous actions (DDPG, TD3,
SAC): it concatenates state and action and returns $Q(s, a)$ with shape
\code{(B, 1)}. \code{DuelingQNetwork} (Wang et al., 2016) combines a value
stream $V(s)$ and an advantage stream $A(s, a)$ as
\begin{equation}
  Q(s, a) = V(s) + A(s, a) - \frac{1}{|\mathcal{A}|}\sum_{a'} A(s, a').
\end{equation}

\begin{pycode}
from toolkit.neural_toolkit import DuelingQNetwork, MLPQNetwork

q = DuelingQNetwork(state_dim=6, action_dim=3, hidden_dims=(64, 64))
states, actions = torch.randn(5, 6), torch.tensor([0, 1, 2, 1, 0])
print(q(states).shape, q(states, actions).shape)   # (5, 3) and (5,)

critic = MLPQNetwork(state_dim=6, action_dim=2, continuous_action=True)
print(critic(states, torch.randn(5, 2)).shape)     # (5, 1)
\end{pycode}

\paragraph{Variational autoencoder.}
\code{VariationalEncoder} returns \code{(z, mu, logvar)} with
$z = \mu + \exp(\tfrac{1}{2}\log\sigma^2)\,\varepsilon$ and
$\varepsilon \sim \mathcal{N}(0, I)$ drawn from PyTorch's default generator
(the reparameterisation trick). The static method \code{kl\_divergence(mu,
logvar)} returns the per-sample KL term of the evidence lower bound,
\begin{equation}
  \mathrm{KL}\big(\mathcal{N}(\mu, \operatorname{diag}\sigma^2)\,\big\|\,\mathcal{N}(0, I)\big)
  = -\frac{1}{2}\sum_{j} \left(1 + \log\sigma_j^2 - \mu_j^2 - \sigma_j^2\right),
\end{equation}
with shape \code{(B,)}. \code{VariationalDecoder} is an MLP from $z$ to the
reconstruction.

\begin{pycode}
import torch.nn.functional as F
from toolkit.neural_toolkit import VariationalDecoder, VariationalEncoder

encoder = VariationalEncoder(input_dim=64, latent_dim=8, hidden_dims=(128,))
decoder = VariationalDecoder(latent_dim=8, output_dim=64, hidden_dims=(128,))
x = torch.randn(32, 64)
z, mu, logvar = encoder(x)
reconstruction = F.mse_loss(decoder(z), x, reduction="none").sum(dim=1)
loss = (reconstruction + encoder.kl_divergence(mu, logvar)).mean()
\end{pycode}

\paragraph{Sequence decoders.}
\code{RNNDecoder} feeds the projected latent vector to a unidirectional LSTM
or GRU at every time step and maps each output with a shared MLP head.
\code{TransformerDecoder} is non-autoregressive: the queries are positional
encodings of the output positions, the memory is the projected latent vector,
and all positions are decoded in parallel. Both generate \code{max\_seq\_len}
steps unless \code{forward(z, seq\_len=...)} asks for another length.

\section{Factories}\label{sec:neural-factories}

Five factories build networks from a type name, which is convenient when the
architecture is a command-line option. Names are case-insensitive; an unknown
name raises a \code{ValueError} that lists the valid names. The remaining
keyword arguments are passed to the constructor of the selected class.

\begin{center}
  \small
  \begin{tabular}{@{}lll@{}}
    \toprule
    Factory & Static method & Valid type names \\
    \midrule
    \code{PolicyFactory}   & \code{create\_policy}         & \code{mlp}, \code{cnn}, \code{rnn}, \code{transformer} \\
    \code{ValueFactory}    & \code{create\_value\_network} & \code{mlp}, \code{cnn}, \code{rnn}, \code{transformer} \\
    \code{QNetworkFactory} & \code{create\_q\_network}     & \code{mlp}, \code{dueling}, \code{cnn}, \code{rnn}, \code{transformer} \\
    \code{EncoderFactory}  & \code{create\_encoder}        & \code{mlp}, \code{cnn}, \code{rnn}, \code{transformer}, \code{vae} \\
    \code{DecoderFactory}  & \code{create\_decoder}        & \code{mlp}, \code{cnn}, \code{rnn}, \code{transformer}, \code{vae} \\
    \bottomrule
  \end{tabular}
\end{center}

Each method takes the type name as its first argument followed by the keyword
arguments of the constructor.

\begin{pycode}
from toolkit.neural_toolkit import PolicyFactory, QNetworkFactory

policy = PolicyFactory.create_policy("mlp", input_dim=8, output_dim=4,
                                    hidden_dims=(64, 64))
q_net = QNetworkFactory.create_q_network("dueling", state_dim=8, action_dim=4)
\end{pycode}

\section{NetworkUtils}\label{sec:neural-utils}

\code{NetworkUtils} groups helper functions as static methods
(Table~\ref{tab:neural-utils}).

\begin{table}[htbp]
  \centering
  \small
  \caption{Static methods of \code{NetworkUtils}.}
  \label{tab:neural-utils}
  \renewcommand{\arraystretch}{1.05}%
  \begin{tabularx}{\textwidth}{@{}>{\raggedright\arraybackslash}p{6.5cm}>{\raggedright\arraybackslash}X@{}}
    \toprule
    Method & Description \\
    \midrule
    \multicolumn{2}{@{}l}{\emph{Initialisation and parameters}} \\
    \code{initialize\_weights(module, method="xavier\_uniform", gain=1.0)} &
      Initialise every linear and (transposed) convolution layer recursively;
      biases are set to zero. Methods: \code{xavier\_uniform},
      \code{xavier\_normal}, \code{kaiming\_uniform}, \code{kaiming\_normal},
      \code{orthogonal}, \code{uniform} ($U(-0.1, 0.1)$), \code{normal}
      ($\mathcal{N}(0, 0.1^2)$). \\
    \code{count\_parameters(model)} & Number of trainable parameters. \\
    \code{count\_parameters\_by\_layer(model)} & Trainable parameters owned directly by each sub-module. \\
    \code{freeze\_layers(model, names)}, \code{unfreeze\_layers(model, names)} &
      Disable or enable gradients of parameters whose name contains any of
      \code{names}. \\
    \midrule
    \multicolumn{2}{@{}l}{\emph{Gradients}} \\
    \code{get\_grad\_norm(model, norm\_type=2.0)} & Total gradient norm. \\
    \code{clip\_grad\_norm(model, max\_norm, norm\_type=2.0)} & Clip in place; returns the norm before clipping. \\
    \midrule
    \multicolumn{2}{@{}l}{\emph{Optimisation}} \\
    \code{create\_optimizer(model, optimizer\_type="adam", lr=1e-3, weight\_decay=0.0, skip\_list=None, **kw)} &
      \code{adam}, \code{adamw}, \code{sgd}, \code{rmsprop}, \code{adagrad},
      \code{adamax}. With \code{weight\_decay > 0} the decay is applied to
      weights only (see \code{apply\_weight\_decay}). \\
    \code{apply\_weight\_decay(model, weight\_decay, skip\_list=None)} &
      Parameter groups without decay for one-dimensional parameters (biases,
      normalisation scales) and for names in \code{skip\_list}. \\
    \code{create\_scheduler(optimizer, scheduler\_type="step", **kw)} &
      \code{step}, \code{multistep}, \code{exponential}, \code{cosine},
      \code{cosine\_warm\_restart}, \code{plateau}, \code{linear},
      \code{polynomial}; \code{**kw} go to the PyTorch scheduler (for example
      \code{step\_size} or \code{T\_max}). \\
    \midrule
    \multicolumn{2}{@{}l}{\emph{Checkpoints}} \\
    \code{save\_checkpoint(model, optimizer, epoch, loss, filepath, **extra)} &
      Save model and optimizer state (\code{optimizer} may be \code{None}),
      epoch, loss and any extra entries; atomic write; returns the path. \\
    \code{load\_checkpoint(model, optimizer, filepath, map\_location=None, weights\_only=True, strict=True)} &
      Restore a checkpoint and return \code{(epoch, loss)}. Tensors are loaded
      onto the device of the model by default, so GPU checkpoints load on
      CPU-only machines. \\
    \midrule
    \multicolumn{2}{@{}l}{\emph{Layer construction}} \\
    \code{create\_mlp\_layers}, \code{create\_conv\_layers}, \code{create\_transposed\_conv\_layers} &
      MLP and (transposed) convolution stacks following the conventions of
      Section~\ref{sec:neural-conventions}. \\
    \code{create\_attention\_layer(input\_dim, num\_heads=8, dropout=0.1)} & Batch-first \code{nn.MultiheadAttention}. \\
    \code{create\_positional\_encoding(d\_model, max\_len)} & Frozen sinusoidal table of shape \code{(1, max\_len, d\_model)}; \code{max\_len=5000}. \\
    \code{get\_activation(activation)} & Same as the module-level function. \\
    \midrule
    \multicolumn{2}{@{}l}{\emph{Shape arithmetic}} \\
    \code{compute\_output\_size(n, k, s=1, p=0, d=1)} & $\lfloor (n + 2p - d(k-1) - 1)/s \rfloor + 1$ \\
    \code{compute\_transposed\_output\_size(n, k, s=1, p=0, op=0, d=1)} & $(n-1)s - 2p + d(k-1) + op + 1$ \\
    \code{compute\_conv\_output\_size((H, W), layers)} & Spatial size after a list of layer dictionaries. \\
    \bottomrule
  \end{tabularx}
\end{table}

\section{Tabular tools}\label{sec:neural-tabular}

The tabular tools work on finite state and action spaces with NumPy arrays
and do not need a GPU. Every table owns a \code{numpy.random.Generator}
(argument \code{rng}: a generator, an integer seed, or \code{None} for fresh
entropy) that is used for all sampling, so results are reproducible without
touching the global random state. States and actions are integer indices;
out-of-range indices raise \code{IndexError}.

The three table classes share one constructor:

% norun
\begin{pycode}
QTable(state_space_size, action_space_size=None, initial_value=0.0,
       dtype=np.float32, rng=None)          # also ValueTable, PolicyTable
\end{pycode}

\begin{description}[leftmargin=0.6cm,style=nextline]
  \setlength{\emergencystretch}{3em}% long identifiers; local to this list
  \item[\code{QTable}] Action values $Q[s, a]$ (\code{action\_space\_size}
    is required). Entries are read and written with \code{get\_value},
    \code{set\_value}, \code{update\_value} and \code{get\_q\_values}.
    \code{get\_max\_q\_value} and \code{get\_max\_action} return the
    greedy value and action (the lowest index among ties). \code{get\_policy}
    samples an epsilon-greedy action, and the methods
    \code{get\_epsilon\_greedy\_probs}, \code{get\_softmax\_probs} and
    \code{get\_softmax\_policy} give the corresponding probabilities or a
    softmax sample.
  \item[\code{ValueTable}] State values $V[s]$, with \code{get\_values} for
    the whole array.
  \item[\code{PolicyTable}] A stochastic policy $\pi[s, a]$ that starts
    uniform. \code{set\_value} and \code{update\_value} change one
    probability and rescale the others so that the row still sums to one;
    \code{set\_policy\_probs}, \code{set\_deterministic\_policy},
    \code{get\_policy\_probs} and \code{get\_policy} (sampling) complete
    the interface.
  \item[\code{DiscreteEnvironment}] An abstract base class for small finite
    environments, constructed with \code{(state\_space\_size,
    action\_space\_size, rng=None)}. Subclasses implement \code{reset()},
    which returns the initial state, and \code{step(action)}, which returns
    \code{(state, reward, done, info)}.
\end{description}

\code{update\_value(state, value, action, learning\_rate)} moves an entry
towards a target: $x \leftarrow x + \alpha\,(\text{value} - x)$. The static
methods of \code{DiscreteTools} implement the classic updates on top of it,
with discount $\gamma$ (\code{gamma=0.99}) and step size $\alpha$
(\code{alpha=0.1}). On a terminal transition (\code{done=True}) the bootstrap
term is dropped.
\begin{align}
  \text{Q-learning:}\quad & Q(s,a) \leftarrow Q(s,a) + \alpha\big[r + \gamma \max_{a'} Q(s',a') - Q(s,a)\big], \\
  \text{SARSA:}\quad & Q(s,a) \leftarrow Q(s,a) + \alpha\big[r + \gamma\, Q(s',a') - Q(s,a)\big], \\
  \text{Expected SARSA:}\quad & Q(s,a) \leftarrow Q(s,a) + \alpha\Big[r + \gamma \sum_{a'} \pi(a' \mid s')\, Q(s',a') - Q(s,a)\Big].
\end{align}
For expected SARSA, $\pi$ is a given \code{PolicyTable} or, if none is given,
the epsilon-greedy policy of the Q-table with the given \code{epsilon}.
\code{value\_iteration\_update(value\_table, state, q\_table)} sets
$V(s) = \max_a Q(s,a)$ (the Q-table must already hold one-step look-ahead
values), and \code{policy\_iteration\_update(policy\_table, state, q\_table)}
makes $\pi(\cdot \mid s)$ greedy with respect to $Q$.

The exploration rules are static methods of \code{DiscreteTools} whose first
two arguments are the Q-table and the state; they sample with the Q-table's
generator.
\begin{description}[leftmargin=0.6cm,style=nextline]
  \item[\code{epsilon\_greedy\_policy(q, s, epsilon)}] A uniformly random
    action with probability $\varepsilon$, otherwise the greedy one. The
    corresponding probabilities (with ties sharing the greedy mass) are
    $\pi(a \mid s) = \varepsilon/|\mathcal{A}| + (1-\varepsilon)\,
    \mathbf{1}[a \in G_s]/|G_s|$, where $G_s$ is the set of greedy actions.
  \item[\code{softmax\_policy(q, s, temperature)}] Also available as
    \code{boltzmann\_policy}. Samples from
    $\pi(a \mid s) \propto \exp(Q(s,a)/\tau)$, computed in a numerically
    stable way.
  \item[\code{ucb\_policy(q, s, visit\_counts, exploration\_constant=1.0)}]
    UCB1: $\arg\max_a Q(s,a) + c\sqrt{\ln N(s) / N(s,a)}$ with
    $N(s) = \sum_a N(s,a)$; untried actions are chosen first, at random among
    themselves. \code{visit\_counts} has the shape of the Q-table.
  \item[\code{thompson\_sampling\_policy(q, s, visit\_counts, ...)}] A
    heuristic Beta-Bernoulli Thompson sampler with
    $p = \operatorname{sigmoid}(Q(s,a))$ and posterior
    $\mathrm{Beta}(\alpha_0 + n p,\ \beta_0 + n(1-p))$, $n = N(s,a)$, where
    $\alpha_0$ and $\beta_0$ are \code{prior\_alpha} and \code{prior\_beta}
    (default 1). It is exact only for Bernoulli rewards with $Q$ in logit
    space.
\end{description}

The following example learns to walk right on a six-state chain. The reward
is sparse, so the Q-table starts with optimistic values of 1, which makes the
greedy policy try every action before settling:

\begin{pycode}
import numpy as np
from toolkit.neural_toolkit import DiscreteEnvironment, DiscreteTools

class Chain(DiscreteEnvironment):
    """Chain 0..n-1: action 1 moves right, 0 moves left; reward 1 at the end."""

    def reset(self):
        self.current_state = 0
        return self.current_state

    def step(self, action):
        move = 1 if action == 1 else -1
        last = self.state_space_size - 1
        self.current_state = int(np.clip(self.current_state + move, 0, last))
        done = self.current_state == last
        return self.current_state, float(done), done, {}

env = Chain(state_space_size=6, action_space_size=2, rng=0)
q = DiscreteTools.create_q_table(6, 2, initial_value=1.0, rng=0)   # optimistic
for episode in range(200):
    state = env.reset()
    for t in range(50):
        action = DiscreteTools.epsilon_greedy_policy(q, state, epsilon=0.1)
        next_state, reward, done, _ = env.step(action)
        DiscreteTools.q_learning_update(q, state, action, reward, next_state,
                                        gamma=0.9, alpha=0.5, done=done)
        state = next_state
        if done:
            break
print([q.get_max_action(s) for s in range(5)])    # [1, 1, 1, 1, 1]
\end{pycode}

\section{Example: a DQN agent for LineMsg}\label{sec:neural-example}

The following script trains a small deep Q-network on the LineMsg environment
of Chapter~\ref{chap:networked} with four agents, treating the $2^4$ joint
actions as one discrete action. It shows the pieces that a training loop
typically needs: action-value selection with \code{q\_net(s, a)}, a target
network, an optimizer created by name, gradient clipping and checkpointing.
On a laptop CPU it runs in a few seconds and the greedy policy reaches the
return of the always-relay strategy (about 65 per episode, compared with
about 30 for random actions).

% fresh
\begin{pycode}
from collections import deque

import numpy as np
import torch
import torch.nn.functional as F

import env_lib
from toolkit.neural_toolkit import MLPQNetwork, NetworkUtils

torch.manual_seed(0)
rng = np.random.default_rng(0)
env = env_lib.LineMsgEnv(num_agents=4, max_iter=50)
obs_dim = int(np.prod(env.observation_space.shape))    # 4 agents x 3 cells
n_actions = int(env.action_space.n)                     # 2**4 joint actions

q_net = MLPQNetwork(obs_dim, n_actions, hidden_dims=(64, 64))
target_net = MLPQNetwork(obs_dim, n_actions, hidden_dims=(64, 64))
target_net.load_state_dict(q_net.state_dict())
optimizer = NetworkUtils.create_optimizer(q_net, "adam", lr=1e-3)
buffer = deque(maxlen=10_000)

obs, info = env.reset(seed=0)
for step in range(4000):
    epsilon = max(0.05, 1.0 - step / 2000)
    if rng.random() < epsilon:
        action = int(rng.integers(n_actions))
    else:
        with torch.no_grad():
            action = int(q_net(torch.as_tensor(obs.reshape(1, -1))).argmax())
    next_obs, reward, terminated, truncated, info = env.step(action)
    buffer.append((obs.ravel(), action, reward, next_obs.ravel(), terminated))
    obs = next_obs
    if terminated or truncated:
        obs, info = env.reset()

    if len(buffer) >= 500:
        batch = [buffer[i] for i in rng.integers(len(buffer), size=64)]
        s, a, r, s2, done = (torch.as_tensor(np.array(x)) for x in zip(*batch))
        with torch.no_grad():
            next_q = target_net(s2).max(dim=1).values
            target = r.float() + 0.99 * (1.0 - done.float()) * next_q
        loss = F.smooth_l1_loss(q_net(s, a), target)   # q_net(s, a): Q(s, a), (B,)
        optimizer.zero_grad()
        loss.backward()
        NetworkUtils.clip_grad_norm(q_net, max_norm=10.0)
        optimizer.step()
    if step % 500 == 0:
        target_net.load_state_dict(q_net.state_dict())

NetworkUtils.save_checkpoint(q_net, optimizer, epoch=step, loss=loss.item(),
                             filepath="checkpoints/dqn_linemsg.pt")
\end{pycode}

Evaluating the greedy policy and restoring the checkpoint later:

\begin{pycode}
def evaluate(policy, episodes=10):
    returns = []
    for episode in range(episodes):
        obs, info = env.reset(seed=100 + episode)
        total, done = 0.0, False
        while not done:
            obs, reward, terminated, truncated, info = env.step(policy(obs))
            total += reward
            done = terminated or truncated
        returns.append(total)
    return float(np.mean(returns))

restored = MLPQNetwork(obs_dim, n_actions, hidden_dims=(64, 64))
epoch, last_loss = NetworkUtils.load_checkpoint(restored, None,
                                                "checkpoints/dqn_linemsg.pt")

def greedy(obs):
    with torch.no_grad():
        return int(restored(torch.as_tensor(obs.reshape(1, -1))).argmax())

greedy_return = evaluate(greedy)
random_return = evaluate(lambda o: int(rng.integers(n_actions)))
print(f"greedy {greedy_return:.1f}, random {random_return:.1f}")
\end{pycode}

Two larger examples ship with the repository, in \code{examples/toolkit/}
(Chapter~\ref{chap:examples}):
\begin{itemize}
  \item \code{simple\_transformer\_example.py} trains a
    \code{TransformerPolicyNetwork} with REINFORCE on a memory task;
  \item \code{transformer\_rl\_example.py} trains transformer actor and critic
    networks with PPO on batched environments.
\end{itemize}

\section{Changes in version 1.0}\label{sec:neural-changes}

Version 1.0 fixed several defects that changed the computed functions. The
most important ones are:

\begin{itemize}
  \item The MLP stacks of earlier versions applied no activation between
    hidden layers, so deep ``MLPs'' were nearly linear. This affected every
    policy, value and Q-network, encoder and decoder.
  \item The positional encoding was indexed by batch position instead of
    sequence position.
  \item Convolutional networks with activations other than ReLU crashed; the
    dueling head was combined incorrectly; Q-networks ignored the
    \code{action} argument; \code{CNNEncoder} failed on $64 \times 64$ inputs;
    \code{CNNDecoder} indexed its kernel sizes incorrectly; bidirectional RNNs
    read a backward state that had seen only one time step.
  \item Unknown activation names now raise instead of falling back silently,
    and the \code{device} argument is honoured.
  \item Tabular tools use a per-table random generator instead of the global
    one; UCB, softmax and \code{PolicyTable.set\_value} were corrected.
\end{itemize}

\paragraph{Checkpoint compatibility.}
Because the layer structure changed (activation modules now sit between the
linear layers, so the indices of \code{nn.Sequential} entries differ), model
checkpoints written by versions before 1.0 cannot be loaded into 1.0 networks:
\code{load\_state\_dict} reports missing and unexpected keys. Even where the
parameter shapes happen to match, the networks compute different functions.
Retrain models with version 1.0; Appendix~\ref{app:migration} lists the other
API changes.

\par\endgroup
